\chapter{A Small LISP on Bare Metal}
\label{ch:lisp}\index{LISP}\index{interpreter}

The earlier chapters built drivers, a filesystem, a network stack and a TLS
implementation -- all in Ada. This one builds a \emph{language}: a small LISP
interpreter, in pure Ada, running an interactive read-eval-print loop on the
ESP32-S3 and able to drive the peripherals from the prompt. It is a capstone that
exercises almost everything underneath -- the heap (Chapter~\ref{ch:heap}), the
octal PSRAM (Chapter~\ref{ch:psram}), exceptions, recursion -- and it turns the
board into a live, scriptable instrument.

LISP is the right choice for this. Its syntax is parentheses, so the reader is a
page; its evaluator is the classic \texttt{eval}/\texttt{apply}, a few hundred
lines; and -- crucially -- it fits the grain of this project. A LISP \emph{closure}
is a data structure, not an Ada closure, so the \texttt{No\_Implicit\_Dynamic\_Code}
restriction that governs every callback in this book never even applies to the
interpreted language; and the bridge to the drivers is Ada calling Ada, with no
foreign-function boundary at all.

\section{One value type}

Everything in the interpreter is an \texttt{Object}: a variant record sized to its
largest case, so every object is the same size and lives as a uniform cell in an
arena.

\begin{lstlisting}
type Object (K : Kind := K_Nil) is record
   Mark : Boolean := False;                -- GC mark bit
   case K is
      when K_Cons    => Car, Cdr : Ref;    -- a pair
      when K_Int     => I : Long_Long_Integer;
      when K_Bool    => B : Boolean;       -- #t / #f
      when K_Symbol  => Sym : Symbol_Id;   -- interned
      when K_Closure => Params, Code, Env : Ref;
      when K_Prim    => Fn : Prim_Fn; Fn_Name : Symbol_Id;
      when K_Nil     => null;              -- the empty list ()
   end case;
end record;
\end{lstlisting}

A \texttt{Ref} is an access into the arena. Symbols are \emph{interned} in a side
table so that two reads of the same name yield the same object (LISP \texttt{eq?}
on symbols is pointer equality), and \texttt{nil}/\texttt{\#t}/\texttt{\#f} are
canonical singletons. The reader turns text into these objects; the printer turns
them back. None of it needs hardware, so the whole core is host-tested natively
before it ever reaches the board.

\section{Putting the arena in PSRAM}

The arena is the one part that the chip's memory map dictates. A LISP wants tens of
thousands of cells; at sixteen-odd bytes each that is hundreds of kilobytes, far
more than the internal SRAM holds. The first build proved the point by overflowing
the data segment. The fix is to allocate the arena from the \emph{heap}, which the
build places in the 8\,MB PSRAM:

\begin{lstlisting}
Arena := new Cell_Array (1 .. Cells);     -- lands in the PSRAM heap
\end{lstlisting}

That raised a second, subtler issue: building the global environment allocates
cells, and if it runs during \emph{elaboration} -- before the heap is up -- it
fails. So the interpreter does no allocation at elaboration; \texttt{Lisp.Init} and
\texttt{Lisp.Eval.Init} are explicit calls the application makes from \texttt{Main},
once PSRAM and the heap are ready.

\section{The console is the USB-serial-JTAG}

Output already goes to the USB-serial console through the ROM. Input comes from the
same port, read at the register level: the \texttt{USB\_DEVICE} peripheral's EP1
endpoint exposes a byte FIFO and a ``data available'' flag.

\begin{lstlisting}
function Try_Get (C : out Character) return Boolean is
begin
   if USB_DEVICE_Periph.EP1_CONF.SERIAL_OUT_EP_DATA_AVAIL then
      C := Character'Val (Natural (USB_DEVICE_Periph.EP1.RDWR_BYTE));
      return True;
   end if;
   return False;
end Try_Get;
\end{lstlisting}

A line reader accumulates characters until a newline, echoing as it goes, and the
REPL is then the obvious loop: read a line, \texttt{Read} it to an object,
\texttt{Eval} it, \texttt{Print} the result, catch \texttt{Lisp\_Error} so a bad
form never stops the loop. On the board:

\begin{lstlisting}[style=sh]
lisp> (define (sq n) (* n n))            sq
lisp> (sq 12)                            144
lisp> (define (fact n) (if (< n 2) 1 (* n (fact (- n 1)))))
lisp> (fact 10)                          3628800
\end{lstlisting}

\section{Garbage collection, at the one safe point}

Without reclamation the arena only grows, so a long session eventually exhausts it.
The collector is a textbook mark-sweep, with one design decision that makes it both
correct and simple: it runs \emph{only between top-level forms}. At the prompt,
after a form has been evaluated and printed, the only live objects are those
reachable from the global environment -- nothing is held half-computed in an Ada
local. So the global environment is the single root: mark everything reachable from
it, free the rest.

\begin{lstlisting}
function GC (Root : Ref) return Natural;   -- mark from Root, sweep, return reclaimed
\end{lstlisting}

Definitions survive because they live in the global environment; a form's garbage
-- its parse tree, its environment frames, the intermediate values, the result just
printed -- does not, and is reclaimed. The REPL calls \texttt{GC} after each form,
and \texttt{(room)} reports the cells in use: it stays flat across forms where it
would otherwise climb without bound.

\cnote{the cells are \texttt{aliased} (so \texttt{'Access} yields a \texttt{Ref}),
and an aliased object is \emph{constrained by its discriminant} -- you cannot change
a cell's \texttt{Kind} through \texttt{Ref.all}. So the free list is a separate array
of \emph{indices}, a cell is rewritten by direct array assignment (\texttt{Arena (I)
:= Template}, which does permit the kind to change), and the sweep only relinks
indices -- it never rewrites a cell. A small rule with a real consequence for the
data structure.}

\section{Tail calls}

\texttt{eval} is naturally recursive, and a bare-metal task stack is small, so deep
recursion is a hazard. The fix is tail-call optimisation: \texttt{eval} is written
as a loop, and the \emph{tail} positions -- the taken branch of an \texttt{if}, the
last form of a \texttt{begin}, and the last body form of a function call -- update
the current expression and environment and go round the loop instead of recursing.

\begin{lstlisting}
when K_Closure =>
   ... bind the parameters in New_Env ...
   Cur_Env  := New_Env;
   Cur_Expr := <last body form>;          -- a tail call: loop, no stack growth
\end{lstlisting}

A call is in \emph{tail} position when its result is the result of the enclosing
function -- there is nothing left to do once it returns. Counting down is
tail-recursive; the recursive call is the whole of the \texttt{if} branch:

\begin{lstlisting}
(define (count n)
  (if (= n 0) 'done
      (count (- n 1))))           ; tail: its value IS the branch's value
\end{lstlisting}

Factorial, written the obvious way, is \emph{not}: the recursive call has to return
before the multiply can run, so every call must keep a frame and the depth becomes
the stack.

\begin{lstlisting}
(define (fact n)
  (if (< n 2) 1
      (* n (fact (- n 1)))))      ; NOT tail: the * is waiting on fact
\end{lstlisting}

The standard way to make an accumulating loop tail-recursive is to carry the running
result in an extra argument, so the recursive call is once more the last thing done:

\begin{lstlisting}
(define (sum n acc)
  (if (= n 0) acc
      (sum (- n 1) (+ acc n))))   ; tail: the + is in the argument, before the call
\end{lstlisting}

Such a loop runs in constant Ada stack. On the board,
\texttt{(count 3000)} on a 3000-deep tail recursion returns \texttt{done} where
without this it would overflow the 64\,KB task stack and crash. The remaining bound
is the arena, not the stack -- \texttt{(count 5000)} reports a clean
\texttt{out of memory} (there is no mid-form GC yet) and the REPL carries on.

\section{Reaching the drivers}

Here is the payoff. Because the interpreter and the HAL are one Ada binary, exposing
a driver is not foreign-function glue -- it is an Ada function, registered as a LISP
primitive, that calls the HAL directly:

\begin{lstlisting}
function Prim_Gpio_Out (Args : Ref) return Ref is
   Pin : constant GPIO.Pin_Id := GPIO.Pin_Id (Int_Value (Car (Args)));
begin
   GPIO.Configure (Pin, GPIO.Output);
   GPIO.Write (Pin, Is_Truthy (Car (Cdr (Args))));
   return Car (Cdr (Args));
end Prim_Gpio_Out;
...
Lisp.Eval.Register_Primitive ("gpio-out", Prim_Gpio_Out'Access);
\end{lstlisting}

Each primitive is library-level and closure-free -- the same rule the HAL callbacks
obey -- so \texttt{'Access} of it is well-formed. The interpreter core
(\texttt{libs/lisp}) stays free of any HAL dependency, so it remains host-testable;
the bridge lives in the application. With GPIO and ADC bound, the board answers:

\begin{lstlisting}[style=sh]
lisp> (gpio-out 2 #t)                     #t      ; LED on
lisp> (adc-read 3)                        2887    ; a real 12-bit sample
lisp> (define (blink n) (if (> n 0) (begin (gpio-toggle 2) (blink (- n 1))) 'done))
lisp> (blink 8)                           done    ; a LISP loop driving a pin
\end{lstlisting}

\subsection{Stateful drivers: a handle registry}

GPIO and ADC are addressed by number, but a session-based driver -- SPI, the W5500,
an ext4 mount, a Modbus connection -- needs the interpreter to refer back to a
specific object. Storing a raw Ada access to that object inside a GC-managed LISP
value would invite exactly the accessibility and lifetime traps of
Chapter~\ref{ch:debug}. The clean answer is a \emph{handle registry}: keep the
sessions in a library-level table, and let the LISP handle be a small integer index
into it.

\begin{lstlisting}
Sessions : array (1 .. Max_SPI) of aliased SPI.Session;   -- library-level
-- (spi-open) acquires Sessions (N) and returns N as an integer handle;
-- (spi-xfer HANDLE BYTES) looks the session up and transfers.
\end{lstlisting}

The session outlives the LISP handle and no access escapes into the heap, so no
accessibility rule is bent. Bytes pass as ordinary LISP lists, converted to a stack
buffer (internal SRAM, so the DMA can reach it) at the boundary. The result, read
from the on-board flash at the prompt, is the chip answering for itself:

\begin{lstlisting}[style=sh]
lisp> (define s (spi-open))
lisp> (spi-xfer s (list #x9f 0 0 0))      (255 239 64 25)
\end{lstlisting}

\texttt{EF 40 19} is the W25Q256 JEDEC identity -- Winbond, 256\,Mbit -- read over SPI
from a language that did not exist when the board powered on. The same pattern
extends to every stateful driver in the HAL, which is the point: a REPL that reaches
the whole stack is the embedded equivalent of a Python prompt on a workstation, but
pure Ada, on bare metal, with a garbage collector you can read in an afternoon.

\noindent\textbf{Example:} \texttt{examples/esp32s3\_lisp} -- the REPL, with the
GPIO/ADC/SPI bridge. The interpreter is \texttt{libs/lisp} (host-tested 53/53:
reader, evaluator, closures, recursion, GC, tail calls).
